\emph{(i)} Let $B$ be invariant. For a bivector $b$, $\rev{b} = -b$, so invariance reads $B(b\psi, \phi) = B(\psi, \rev{b}\phi)$. If the identity $B(a\psi, \phi) = B(\psi, \rev{a}\phi)$ holds for $a$ and for $a'$, it holds for $aa'$: $B(aa'\psi, \phi) = B(a'\psi, \rev{a}\phi) = B(\psi, \rev{a'}\rev{a}\phi) = B(\psi, \rev{(aa')}\phi)$; it holds for scalars and is linear in $a$. Every even blade is a product of bivector blades, $e_{i_{1}}e_{i_{2}}\cdots e_{i_{2k-1}}e_{i_{2k}} = (e_{i_{1}}e_{i_{2}})\cdots(e_{i_{2k-1}}e_{i_{2k}})$, so the identity holds for every $a \in \Cl^{0}(V)$ (the check confirms for $2 \le n \le 7$, all signatures, that products of bivector blades reach every even blade). Every element of $S$ is $\psi = af$ for some $a \in \Cl^{0}(V)$, and $f = ff$. Hence, for $\phi \in S'$,
\[
B(af, \phi) = B(f, \rev{a}\phi) = B(ff, \rev{a}\phi) = B\bigl(f, \rev{f}\rev{a}\phi\bigr) = B\bigl(f, \rev{(af)}\phi\bigr) = B(f, \rev{\psi}\phi),
\]
where $\rev{\psi}\phi = \rev{f}\rev{a}\phi$ lies in $\Cl^{0}(V)f' = S'$, the domain of $B(f, \cdot)$. Conversely, for any linear $\lambda$ on $\Cl^{0}(V)$ and bivector $b$, $\lambda(\rev{(b\psi)}\phi + \rev{\psi}b\phi) = \lambda(\rev{\psi}(\rev{b} + b)\phi) = 0$. So the invariant forms are exactly the linear functionals of the reversion pairing, and they all vanish if and only if the pairing does.

\emph{(ii)} By Proposition \claimref{synthesis.reversion-pairing-vanishes-residues-zero-four}(i), at a causal signature the conditions $n$ even, $\omega^{2} = +1$, $\rev{\omega} = -\omega$ hold exactly at residues $0$ and $4$, where $n \equiv 2 \pmod 4$; there $P_{\pm} = \tfrac12(1 \pm \omega)$ are central idempotents of $\Cl^{0}(V)$ with $\rev{P_{+}} = P_{-}$. For $\psi, \phi$ in ideals of chirality $+$ (in $\Cl^{0}(V)P_{+}$, any number of copies), $\rev{\psi}\phi = \rev{\psi}P_{+}\phi$ and $\rev{\psi} = \rev{(\psi P_{+})} = P_{-}\rev{\psi}$ with $P_{-}$ central, so $\rev{\psi}\phi = \rev{\psi}P_{-}P_{+}\phi = 0$; by (i) every invariant form vanishes. Between the chiralities $\rev{P_{+}}P_{-} = P_{-} \neq 0$ and the cross forms exist (computed in (iii)). At the other causal residues the reversion pairing on $S$ is not identically zero (Proposition \claimref{synthesis.reversion-pairing-vanishes-residues-zero-four}(ii)), and each of its blade components is an invariant form by the converse in (i). A mass term of a field valued in $S$ (or in several copies of $S$, or a complex combination, read as a real form on the underlying real space) is a derivative-free bilinear expression invariant under the connected spin group, hence under its Lie algebra; for anticommuting fields it uses an alternating form, for a pair of fields any form. Every such form is invariant in the sense above, so none is nonzero at residues $0$ and $4$ for one ideal. The dimension-six ideal is that of Proposition \claimref{synthesis.lowest-causal-residues-zero-four-content}, with the chiral factor in $f$.

\emph{(iii)} The check builds $S$ at each signature as in Definition \claimref{obj.dark-tower-matter.ideals-fields-momentum}, writes the left action of each basis bivector as an exact matrix on a basis of $S$, and solves the linear conditions $B(b\psi, \phi) + B(\psi, b\phi) = 0$ for the matrix of $B$; the kernel dimensions are those stated. It then forms, for every blade $C$, the bilinear form $\psi, \phi \mapsto$ (coefficient of $C$ in $\rev{\psi}\phi$), and verifies that these forms span a space of the same dimension and that adjoining the kernel does not raise the rank, which is (i) in coordinates. At $(1,4)$ and $(4,1)$ it also verifies $B(a\psi, \phi) = B(\psi, \rev{a}\phi)$ for every even blade $a$ and every invariant form. The symmetry types are read off the kernel bases.

\emph{(iv)} For one constituent the pairing in question is $B(\psi, \psi)$; over a field of characteristic zero it vanishes for every $\psi$ exactly when $B$ is alternating. So some invariant form is nonzero on a single constituent exactly when some invariant form is not alternating; otherwise a nonzero pairing needs two constituents, $B(\psi_{1}, \psi_{2})$, which exists when some form is nonzero. With (iii): two at $n = 3, 4$, one at $n = 5$, none at $n = 2, 6$. Theorem \claimref{dark-tower-matter.where-mass-begins} proves mass begins at two constituents for $n = 3, 4$ (the reversion pairing is alternating there) and at one for $n = 5$; the thresholds agree. The agreement concerns where an invariant pairing first exists, not the mass formulas themselves, which are the theorem's.

\emph{Scope.} The result is Clifford structure, transport-absent. It does not decide between the pairing and momentum readings of mass in dimensions $n \equiv 2 \pmod 4$; it shows that the pairing reading has no internal freedom (every invariant pairing is the reversion pairing read through a functional), and that it coincides with the field-theoretic notion of a mass term. The standard field-theory statement that a single chirality has no mass term in dimensions $2 \bmod 4$ (cited from memory, for comparison only) is recovered here with an own proof and is not used. The factorisation (i) is in substance the classical description of spin-invariant bilinear forms on spinors through the reversion and conjugation anti-involutions (Harvey, \emph{Spinors and Calibrations}; Alekseevsky and Cort\'es; both cited from memory, for comparison only); the proof given here is self-contained and uses neither. The argument of (i) does not use nondegeneracy of the quadratic form beyond the existence of the bivector action, so it also covers degenerate signatures, which are not claimed here.
